\section*{Chapter \ref{ch:dv:stochastic-volatility} --- Stochastic Volatility}

\begin{solution}{exo:dv:stochastic-volatility:1}
Base: $2\times1.5\times0.04-0.6^2=-0.24$; calibrated: $2\times2.05\times0.060-0.66^2=-0.19$. Both fail: the
variance can reach zero, which the characteristic function handles, but a naive simulation must truncate.
\end{solution}

\begin{solution}{exo:dv:stochastic-volatility:2}
$v(t)=\bar v+(v_0-\bar v)e^{-\kappa t}$: 0.0589 at six months, 0.0489 at one year. The one-year total variance is
$0.04+0.04(1-e^{-1.5})/1.5=0.0607$: a volatility of 24.6\%.
\end{solution}

\begin{solution}{exo:dv:stochastic-volatility:3}
$\varphi(-i)=\E[e^{x_T}]=\E[S_T]/F=1$ because the forward is the expectation of $S_T$ under the pricing measure. A
value of 1.01 means the model's forward is 1\% off: every call is mispriced by about 1\% of the forward, an
arbitrage against the futures market, usually a sign of a wrong branch or a coding error.
\end{solution}

\begin{solution}{exo:dv:stochastic-volatility:4}
Under the pricing measure the drift of the variance already contains it: calibration to options estimates
$\kappa\bar v$ and $\kappa$ as they are priced, not as they behave. Historical volatility estimates the real-world
level, which is typically below the priced one (the variance risk premium of chapter 14); setting $\bar v$ from
history underprices options.
\end{solution}

\begin{solution}{exo:dv:stochastic-volatility:5}
With $\rho=0$ the price is an average of Black--Scholes prices with the forward fixed; in \omterm{def:dv:implied-volatility-and-its-surface:surface}{log-moneyness} the
Black--Scholes implied volatility of such an average is symmetric around $k=0$ (put and call at $\pm k$ correspond by
put--call symmetry). Near the money the Black--Scholes price is nearly linear in volatility, so the average price
corresponds to the average volatility $\E[\bar\sigma]$, which is below $\sqrt{\E[\bar\sigma^2]}$ by Jensen's
inequality.
\end{solution}

\begin{solution}{exo:dv:stochastic-volatility:6}
The short skew is steered by $\rho$ and $\eta$, the long skew too, and $\kappa$ controls how fast the skew decays with
expiry: steepening only the short end needs a higher $\kappa$ and a larger $\eta$ at once, which also changes the
curvature and the term structure. Heston's skew term structure has a fixed shape, and a local event in one part of
it cannot be fitted by moving one parameter.
\end{solution}

\begin{solution}{exo:dv:stochastic-volatility:7}
$\eta=0.77$; the one-year 90--110 skew rises to 6.45 points, since with a negative correlation the \omterm{def:dv:stochastic-volatility:sv}{volatility of
volatility} also steepens the skew.
\end{solution}

\begin{solution}{exo:dv:stochastic-volatility:8}
The model was calibrated where it fits; at one week its skew is half the market's, so it underprices short
out-of-the-money puts and gives the wrong hedges there. Price weeklies from the market's own short-dated smile, or
with a model that produces steep short skews (jumps, rough volatility).
\end{solution}

\begin{solution}{pb:dv:stochastic-volatility:1}
\textbf{1.} 23.4\%, 20.1\%, 16.9\%, 14.4\%, 13.3\%.
\textbf{2.} Skew 5.71 points; curvature 0.30 point.
\textbf{3.} No: $2\kappa\bar v-\eta^2=-0.24$.
\textbf{4.} Averaging over random variance lowers the average volatility (Jensen), and the negative correlation
lowers the at-the-money volatility further: 16.9\% at one year.
\textbf{5.} 17.4\%, still below 20\% for the same reason.
\textbf{6.} 3.94 points.
\textbf{7.} $-1.77$ points.
\textbf{8.} A little: the correlation shifts probability between the tails and the forward is fixed, so the
at-the-money price changes by a second-order amount.
\textbf{9.} Risk reversals, and every product whose value depends on the skew: barriers, digitals, autocallables.
\textbf{10.} That volatility tends to rise when the index falls. A 1\% fall is a shock
$dW^1=-0.01/\sqrt{v}=-0.05$; since $d\sqrt v\approx\tfrac12\eta\,dW^2$ and $\E[dW^2\mid dW^1]=\rho\,dW^1$, the
expected rise of $\sqrt v$ is $\tfrac12\times0.6\times0.7\times0.05=0.0105$: about one volatility point.
\textbf{11.} $\eta=0.77$.
\textbf{12.} It steepens to 6.45 points: with a negative correlation the \omterm{def:dv:stochastic-volatility:sv}{volatility of volatility} adds skew as well
as curvature.
\textbf{13.} Make the correlation less negative until the skew is back at its target.
\textbf{14.} $2\times1.5\times0.04-0.77^2=-0.47$.
\textbf{15.} Wing options, variance swaps' convexity (chapter 14), cliquets and volatility-of-volatility products.
\textbf{16.} Parameters interact (here $\eta$ moves the skew too); moving them one at a time leaves the others
wrong.
\textbf{17.} $\kappa$ and $\eta$ are nearly redundant for a surface up to a few years; fixing $\kappa$ removes a flat
direction of the fit and makes the other parameters stable day to day (chapter 24).
\textbf{18.} Its skew levels off where the market's keeps steepening: half the market's at one week.
\textbf{19.} $-1.77$ volatility points for $+0.2$ in $\rho$; $\eta=0.77$ for a curvature of 0.5 point.
\textbf{20.} The uncertainty of future variance, which raises the price of options whose value is convex in
volatility: the wings.
\end{solution}

\begin{solution}{iq:dv:stochastic-volatility:1}
$dS=\sqrt vS\,dW^1$, $dv=\kappa(\bar v-v)dt+\eta\sqrt v\,dW^2$, $d\langle W^1,W^2\rangle=\rho\,dt$. $v_0$: short-dated
level; $\bar v$: long-dated level; $\kappa$: how fast one becomes the other, and how fast the skew fades; $\eta$: the
smile's curvature (wings); $\rho$: its skew.
\iqlookfor{one sentence per parameter, on the surface.}
\end{solution}

\begin{solution}{iq:dv:stochastic-volatility:2}
Two Brownian motions, one traded risky asset: the variance risk cannot be hedged with the share. Hedge the delta with
the share and the \omterm{def:dv:greeks-and-the-hedging-pnl:greeks}{vega} (variance exposure) with another option, typically of the same expiry; the second option makes
the market complete in the model.
\iqlookfor{the count of risks against instruments.}
\end{solution}

\begin{solution}{iq:dv:stochastic-volatility:3}
Condition on the whole variance path: with zero correlation $\ln S_T$ is normal with variance $\int v\,dt$, so the
conditional price is Black--Scholes at that variance; average over variance paths (the \omterm{prop:dv:stochastic-volatility:mixing}{mixing formula}).
\iqlookfor{conditioning and independence.}
\end{solution}

\begin{solution}{iq:dv:stochastic-volatility:4}
The complex logarithm in the characteristic function jumping branches (the ``little Heston trap''): use the form with
$g=(\xi-d)/(\xi+d)$ and $e^{-dT}$, which stays on the principal branch; also check the integration range and the
damping of the Fourier integral.
\iqlookfor{branch cuts, then numerics.}
\end{solution}

\begin{solution}{iq:dv:stochastic-volatility:5}
Check the inputs first (a bad quote, a stale expiry); then check whether the fit has a flat direction ($\rho$ and $\eta$
trading off); stabilise the calibration (fix $\kappa$, penalise parameter changes, chapter 24) and measure the P\&L effect
of the jump on exotics before accepting it; if the market really moved, the exotics' reserves should reflect it.
\iqlookfor{data, identifiability, stability, P\&L impact.}
\end{solution}

\begin{solution}{iq:dv:stochastic-volatility:6}
Its short-dated smile comes only from the diffusion of variance over a short time, so the at-the-money skew tends to a
constant proportional to $\rho\eta/\sqrt{v_0}$ as the expiry shrinks, while the market's grows. Fixes: add jumps
(chapter 13), which produce skew at short expiries, or make volatility rough (chapter 12), which produces a skew
exploding like a power of the expiry.
\iqlookfor{the short-expiry limit, and the two families of fixes.}
\end{solution}
